设 $X>0$。 设 $V\in C^\infty(0,\infty)$，并且对每个
$C>0$ 及每个整数 $j\ge0$，当
$x\to0$ 时 $(x\partial_x)^jV(x)$ 为 $O_{C,j}(x^C)$，当 $x\to\infty$ 时为 $O_{C,j}(x^{-C})$。 定义
\[
 \widehat V(t)=\int_0^\infty V(x)x^t\,\frac{dx}{x}.
\]
待变换的磨光表达式是完备化和
\[
 \begin{split}
 &\frac1{2\pi i}\int_{(a)}\widehat V(s-1/2)X^{s-1/2}T(s,\Psi)\,ds\\
 &\quad=\sum_{\substack{n\ {\rm sf},\ b\\n,b\equiv1\ (3)\\(nb,\mathfrak E)=1}}
 \frac{\gamma_2(n)\overline{\alpha(nb^3)}\Psi(nb^3)}
      {\sqrt{q_n}\,q_b}
 V\!\left(\frac{q_nq_b^3}{X}\right),\qquad a>1.
 \end{split}
\]
此等式由绝对收敛与 Mellin 反演推得。
特征作用在整个指标 $nb^3$ 上，其取零的值保留；
$n$ 与 $b$ 可以有公共素因子。

反射检验函数定义为
\[
 K=\frac{(2\pi)^4}{27},\qquad
 R(t)=\prod_{\pm}\frac{\Gamma(1+t\pm1/6)}{\Gamma(1-t\pm1/6)},
\]
\[
 V^\sharp(x)=\frac1{2\pi i}\int_{(\sigma)}
       \widehat V(-t)R(t)(Kx)^{-t}\,dt\qquad(\sigma\ge0).
\]
定义 $V^\sharp$ 的积分绝对收敛，且与 $\sigma\ge0$ 无关。

乘积 $T$ 具有整函数延拓。 在每个固定竖直
带形域上它具有多项式增长，其中常数允许依赖于
$L,\Psi_0,\mathcal P,(j_p)$ 及该带形域。 存在一组固定的三重
尖点系数函数，从 $\lambda^{-4}\mathcal O\setminus\{0\}$ 到
$\mathbb C$，由证明中的 theta 展开构造。 同一
三重函数用于算术数据、尺度与检验剖面的所有选择。
对上述每个固定的算术选择，存在一族有限的项，
按 (1) 中的方式编号，独立于
$X$、$V$ 及围道参数 $a>1$。 在每一项中，$d$ 取自
这组三重函数，而 $\vartheta$ 是 $\mathcal O$ 的一个加性特征，
模一个仅依赖于 $L$ 的固定模；它们按扇区的依赖关系
恰如 (3) 所述，其余数据具有性质 (1)--(3)。
对这些项及每个 $a>1$，恒等式为
\begin{equation}\label{eq:reflection}
 \begin{aligned}
 &\frac1{2\pi i}\int_{(a)}\widehat V(s-1/2)X^{s-1/2}T(s,\Psi)\,ds\\
 &\quad=\sum_{\text{terms}}\zeta
 \sum_{0\ne\mu\in\lambda^{-4}\mathcal O}
 \frac{d(\mu)\alpha(\mu)\vartheta(\lambda^4\mu)}{\sqrt{q_\mu}}
 \prod_{p\ {\rm active}}B_p(\lambda^4\mu)
 V^\sharp\!\left(\frac{q_\mu X}{q_c^2}\right).
 \end{aligned}
\end{equation}
每个内和绝对收敛。 这些项具有以下
精确性质。

\begin{enumerate}
\item 一项由 $h_0\bmod L$ 及活跃素数的子集确定。
每个 $j_p\ne0$ 的 $p$ 都是活跃的，而 $j_p=0$ 的素数可以是
活跃的或非活跃的。 记 $r=\prod_{p\ {\rm active}}p$。 则
$c=c_Fr$，其中 $c_F$ 是
$\lambda^2h_0/L$ 的既约分母，附以一个正规化单位。 于是 $c_F$ 遍历一个
固定的有限集，且其素因子整除 $L$。 共有
$O_{\mathfrak L}(2^{|\mathcal P|})$ 项，或在把下一式中
每个 Ramanujan 因子拆成它的两个被加项之后为
$O_{\mathfrak L}(4^{|\mathcal P|})$ 项。

\item 局部列因子为
\[
 B_p(x)=\begin{cases}
 \chi_p(x)^{-j_p-2},&j_p\ne0,4,\\
 q_p^{-1/2}(-1+q_p1_{p\mid x}),&j_p=4,\\
 q_p^{-1/2}\chi_p(x)^{-2},&j_p=0\text{ and }p\text{ is active}.
 \end{cases}
\]
函数 $d$ 是三个固定尖点系数函数之一。
把 $u$ 记作 Eisenstein 单位，它们的支撑与大小满足
\begin{equation}
 \operatorname{supp}d\subset
 \{u\lambda^k n b^3:k\ge-4,\ n,b\equiv1\pmod3,\ n\ {\rm sf}\},
 \qquad |d(\mu)|\le27\,3^{k/6}|b|.
 \label{old-eq:4.3}
\end{equation}
该界是针对被支撑指标所展示的表示而言的。
它允许 $n$ 与 $b$ 有公共素数，并且除
$\lambda$ 处所展示的限制外，对 $L$ 的素数不施加任何排除。
特征 $\vartheta$ 是 $\mathcal O$ 的一个商
模一个仅依赖于 $L$ 的固定模的加性特征。

\item 标量 $\zeta$ 与 $\mu$ 无关，且满足
$|\zeta|\le1/81$。 其关于移动素数的依赖可以精确给出。
对 $m\not\equiv0\pmod6$，置
\[
 \tau_{p,m}^{\pm}=q_p^{-1/2}\sum_{y\bmod p}\chi_p(y)^m e(\pm y/p),
 \qquad |\tau_{p,m}^{\pm}|=1.
\]
对每个活跃的 $p$，由
\[
 \sigma_p=\lambda^2c/p,\qquad
 \epsilon_p=-\big((\lambda^3c/p)\sigma_p\big)^{-1},
\]
定义模 $p$ 的单位，并置
\[
 \omega_{p,j}=\begin{cases}
 \tau_{p,j}^{-}\tau_{p,j+2}^{+}\chi_p(\epsilon_p)^{-j-2},&j\ne0,4,\\
 \tau_{p,4}^{-},&j=4,\\
 -\tau_{p,2}^{+}\chi_p(\epsilon_p)^{-2},&j=0.
 \end{cases}
\]
$\tau$ 的下标按模六读取。 令
$M=\lambda^{12}L^4$，存在一个绝对值为一的数 $\kappa_F$，
仅依赖于 $h_0$ 及 $r$ 模 $M^2$ 的类，使得
\begin{equation}\label{eq:reflection-row-phase}
 \zeta=-\frac i{81}\bar\alpha(c)^2\widehat\phi(h_0)\bar\kappa_F
 \prod_{p\ {\rm inactive}}(1-q_p^{-1})
 \prod_{p\ {\rm active}}\chi_p(\sigma_p)^{-2}\omega_{p,j_p}.
\end{equation}
在 $h_0$ 与 $r\bmod M^2$ 的有限多个类中的每一个类里，$d$
与 $\vartheta$ 都是固定的。 因此在该扇区内它们是公共列
因子。 除所展示的尺度 $q_c^2$ 外，
$B_p$ 之外的所有其他移动素数依赖都包含在标量
$\zeta$ 中，而 $\zeta$ 可以依赖于整个活跃集。

分支约定对具有相同 $L$ 与 $\Psi_0$ 的所有局部素数集
同时成立。 特别地，若把一个好素数 $p$ 以指数零加入
局部集并宣布其为非活跃的，则所得
分支具有相同的 $c_F,c,d,\vartheta,\kappa_F$，在每个
其他活跃素数处的数据相同，并与 $p$ 缺席的分支具有相同的核。
其标量乘以 $1-q_p^{-1}$。 这种相容性包括每个
$h_0$，而不只是模 $L$ 的单位类。
\end{enumerate}

\emph{变换后检验函数的界。}
对 $A\ge0$ 及整数 $B\ge0$，置
\[
 \mathfrak M_{A,B}(V)=
 \sup_{-A\le\eta\le1/4}\int_{\mathbb R}
       (1+|u|)^B|\widehat V(\eta+iu)|\,du.
\]
对每个整数 $j\ge0$，
\begin{equation}\label{eq:reflection-kernel}
 |(x\partial_x)^jV^\sharp(x)|
 \ll_{A,j}\min\{x^{1/4},(1+x)^{-A}\}
       \mathfrak M_{A,\lceil4A\rceil+j+2}(V).
\end{equation}
若 $\chi\in C_c^\infty(0,\infty)$ 固定，$Y>0$，且 $J\ge0$ 为
整数，则
\begin{equation}\label{eq:reflection-annular-fourier}
 \int_{\mathbb R}(1+|u|)^J
 \left|\int_0^\infty\chi(y)V^\sharp(Yy)y^{-iu}\frac{dy}{y}\right|du
 \ll_{A,J,\chi}\min\{Y^{1/4},(1+Y)^{-A}\}
       \mathfrak M_{A,\lceil4A\rceil+J+4}(V).
\end{equation}
这些都是有限光滑半范数。 更精确地，
\begin{equation}\label{eq:reflection-test-seminorm}
 \mathfrak M_{A,B}(V)\ll_{A,B}
 \int_0^\infty(1+x^{-A}+x^{1/4})
       \bigl(|V(x)|+|(x\partial_x)^{B+2}V(x)|\bigr)\frac{dx}{x}.
\end{equation}
把 $V(x)$ 换成 $x^{iu_0}V(x)$ 至多使
$\mathfrak M_{A,B}$ 乘以 $(1+|u_0|)^B$。
\end{proposition}

\begin{proof}
我们首先通过有限多个平移后的 theta 值来表达特征掩码。
随后变换它们的有理尖点，并计算所得的自守因子与局部
Fourier 因子；最后的 Mellin 比较给出反射恒等式
及其变换后检验函数的界。

\emph{固定的 theta 输入。}
设 $\Gamma=\mathrm{SL}_2(\mathcal O)$，
$\Gamma_1(3)=\{g\in\Gamma:g\equiv I\pmod3\}$，且
$\Gamma_2=\langle\mathrm{SL}_2(\mathbb Z),\Gamma_1(3)\rangle$。
\cite[方程 (5.4)、(5.7)、(5.9) 与 (5.12)--(5.17)]{DR}
中无条件成立的公式给出
$\mathbb H^3=\mathbb C\times\mathbb R_{>0}$ 上的光滑三次 theta 函数 $\theta$，满足
\[
 \theta(gw)=\kappa(g)\theta(w)\quad(g\in\Gamma_2),\qquad
 \kappa\!\begin{pmatrix}a&b\\c&\delta'\end{pmatrix}
   =\left(\frac ca\right)_3\quad
   \left(\begin{pmatrix}a&b\\c&\delta'\end{pmatrix}\in\Gamma_1(3),\ c\ne0\right).
\]
这里 $\kappa$ 在 $\mathrm{SL}_2(\mathbb Z)$ 上为一，且在 $\Gamma_1(3)$ 中
左下元素为零的元素上为一。 下标 $3$
表示通常的三次剩余符号；对好素数 $p\notin S$，
有 $(x/p)_3=\chi_p(x)^2$。
这些自守与系数公式不假设任何关于 $L$ 函数零点的
条件。

置
\[
 \mathsf L(t)=\begin{pmatrix}1&0\\t&1\end{pmatrix},\qquad
 (t_0,t_1,t_2)=(0,\omega^2,\omega),\qquad
 \Theta_b(w)=\theta(\mathsf L(t_b)w).
\]
若 $v\in\mathbb Z+3\mathcal O$，则 $\mathsf L(v)$ 是
$\mathrm{SL}_2(\mathbb Z)$ 的一个元素与
$\Gamma_1(3)$ 中一个自守因子为一的下平移之积。 对
上平移有类似断言成立。 因此
\[
 \theta(\mathsf L(\lambda)w)=\Theta_1(w),\qquad
 \theta(\mathsf L(-\lambda)w)=\Theta_2(w),
\]
因为 $\lambda-\omega^2=2+3\omega$ 且
$-\lambda-\omega=-1-3\omega$。 若 $u_0=s+t\omega$，其中
$0\le s,t<3$，且 $E=\bigl(\begin{smallmatrix}0&-1\\1&0\end{smallmatrix}\bigr)$，
则
\[
 \begin{pmatrix}u_0&-1\\1&0\end{pmatrix}=E\mathsf L(-u_0),
 \qquad -u_0-t_t\in\mathbb Z+3\mathcal O,
\]
从而 $\theta\bigl(\bigl(\begin{smallmatrix}u_0&-1\\1&0\end{smallmatrix}\bigr)w\bigr)
=\Theta_t(w)$。 这些是下文所需的仅有的尖点代表元。

对其中任一个，由
\begin{equation}\label{eq:reflection-cusp-expansion}
 \bar\theta(H(z,v))=C_H(v)+
 \sum_{\mu\ne0}d_H(\mu)vK_{1/3}(4\pi|\mu|v)\check e(\mu z).
\end{equation}
定义 $d_H$。 这里 $C_H(v)$ 与 $z$ 无关。
三个函数 $\Theta_0,\Theta_1,\Theta_2$ 分别是
~\cite{DR} 中的 $F_1,F_{19},F_{10}$。 其附录~A 的第 $1,19,10$ 行
用其方程 (5.7)、(5.13) 与 (5.14) 中的 $\tau,\tau_1,\tau_2$
表达它们的系数。 在其记号中，共轭
展开使用 $\overline{d_j(-\mu)}$，如它们的方程 (5.16)。
那些公式给出
方程~\eqref{old-eq:4.3} 中的支撑包含。 为完整起见，$\tau$ 在
$\mu=u\lambda^k n b^3$ 处可能的两个大小为
\[
 3^{k/6+8/3}|b|\quad\hbox{and}\quad3^{k/6+3}|b|,
\]
而 $\tau_1,\tau_2$ 的大小为 $9|b|$ 且 $k=-4$。
这里 Gauss 和仅使用三次符号，包括在 $S$ 中的素数处。
为明确其定义域，对每个主无平方因子 $n$，如
~\cite[方程 (1.7)]{DR} 中那样定义
\begin{equation}\label{eq:reflection-cubic-source-gauss}
 \widetilde g_3(n)=q_n^{-1/2}\sum_{v\bmod n}
       \left(\frac vn\right)_3\check e(v/n),
 \qquad \widetilde g_3(1)=1.
\end{equation}
通常的三次符号在除 $\lambda$ 外的每个素数处有定义，
在非单位上取值为零，并在分母中按乘性延拓。
每个主 $n$ 都与 $\lambda$ 互素。 在 $n$ 的每个素因子处
三次特征非平凡，且 $\check e(v/n)=e(\lambda v/n)$ 是本原的，
因为 $\lambda$ 在那里是单位。 因此有限域 Gauss 恒等式与
中国剩余定理给出 $|\widetilde g_3(n)|=1$，对每个
这样的无平方因子 $n$ 成立。 这包括范数为 $4$ 的主素数 $-2$；
不使用以该数为分母的六次特征。 又，只要 $(t,n)=1$，
变量替换 $y=tv$ 给出
\[
 q_n^{-1/2}\sum_{v\bmod n}\left(\frac vn\right)_3\check e(tv/n)
       =\left(\frac tn\right)_3^{-1}\widetilde g_3(n).
\]
这涵盖了其他源系数族中的扭曲 $t=\lambda^2,\omega\lambda^2,\omega^2\lambda^2$。
于是所有三个尖点族中出现的未正规化 Gauss 和
其大小均为 $|n|$。 从而
上面展示的每个系数大小至多为
$27\,3^{k/6}|b|$，包括 $k=-4$。 这些公式只要求
$n$ 无平方因子且 $n,b$ 为主；它们不要求
$(n,b)=1$。 $nb^3$ 的素赋值还表明它由这样的 $n,b$
的表示（当存在时）是唯一的。 特别地，当
这些变量有公共素数或含有被排除在素数求和之外的素数时，同样的界
仍然成立。

在无穷远处，条件 $x=\lambda^3\mu\equiv1\pmod3$ 恰好分离出
$\mu=\lambda^{-3}nb^3$，其中 $n,b$ 为主且 $n$ 无平方因子。
无穷远处的支撑有 $\lambda^3\mu\in\mathcal O$；~\cite{DR}
方程 (5.7) 中其他每个分歧
指数都使 $x$ 被 $\lambda$ 整除，而剩余
一对中的负号给出 $x\equiv-1\pmod3$。 在被分离出的指标处，$\bar\theta$
中的系数为
\begin{equation}\label{eq:reflection-infinity-coefficient}
 d_I(\lambda^{-3}nb^3)=3^{5/2}|b|\widetilde g_3(n).
\end{equation}
这是未加掩码的源公式，对每个主无平方因子 $n$
与每个主 $b$ 成立，它们之间没有互素条件。 仅当
$(n,S)=1$ 时，它才可以用整体六次记号表达：在 $\gamma_2(n)$ 中作替换
$y=\lambda v$ 给出
\begin{equation}\label{eq:reflection-cubic-sextic-conversion}
 \widetilde g_3(n)=\left(\frac\lambda n\right)_3^{-1}\gamma_2(n)
       =\chi_n(\lambda)^{-2}\gamma_2(n)\qquad((n,S)=1).
\end{equation}
~\cite[方程 (1.5)]{DR} 中的三次补充律使
$(\lambda/x)_3$ 在所有主 $x$ 上模 $9$ 周期。 在与 $S$ 互素的
主元素上它等于 $\chi_x(\lambda)^2$。 $\mathcal X$ 中的特征在与 $\mathfrak E$
互素的主元素上模其导子周期，而
零条件 $(x,\mathfrak E)\ne1$ 模 $\mathfrak E$ 周期。
因此一个具有素支撑 $S$ 的充分可整除的 $L$ 对
固定有限族的每个成员都适用。 这在选择移动局部素数
之前验证了它的选取。

\emph{带零延拓的有限 Fourier 反演。}
对 $j\in\{0,\ldots,5\}$ 定义
\[
 C_{p,j}(h)=q_p^{-1}\sum_{x\bmod p}\chi_p(x)^j e(-hx/p).
\]
由 $e$ 给出的特征配对在
$\mathcal O/(p)$ 上非退化。 在 Gauss 和中作变量替换，并利用
非平凡乘性特征的正交性，给出
\begin{equation}\label{eq:reflection-local-fourier}
 C_{p,j}(h)=\begin{cases}
 q_p^{-1/2}\tau_{p,j}^{-}\chi_p(h)^{-j},&j\ne0,\ h\ne0,\\
 0,&j\ne0,\ h=0,\\
 -q_p^{-1},&j=0,\ h\ne0,\\
 1-q_p^{-1},&j=0,\ h=0.
 \end{cases}
\end{equation}
这里及下文 $h=0$ 表示 $h\equiv0\pmod p$。 $\chi_p$ 的每个非平凡幂
在素数 $p$ 处都是本原的，因此通常的有限域
Gauss 和恒等式给出 $|\tau_{p,m}^{\pm}|=1$。

$\mathcal O/(L)$ 上及每个 $\mathcal O/(p)$ 上的
Fourier 反演现在给出
\[
 \phi(x)\prod_p\chi_p(x)^{j_p}
  =\sum_h C(h)e\!\left(x\left(h_0/L+\sum_p h_p/p\right)\right),
 \qquad C(h)=\widehat\phi(h_0)\prod_pC_{p,j_p}(h_p).
\]
这用到了每个加性频率，包括非单位 $h_0$ 与零局部
频率。 我们不引用
\cite[引理~5.2 与推论~5.1]{DR} 中的扭曲公式，其单位 Fourier 支撑假设
对这些掩码未必成立。 以其无穷远系数 $x=\lambda^3\mu\in\mathcal O$
上带有该自守因子的 theta 函数是有限和
\begin{equation}\label{eq:reflection-translated-theta}
 \mathcal F(z,v)=\sum_hC(h)\bar\theta(z+z_h,v),\qquad
 z_h=\lambda^2\left(h_0/L+\sum_p h_p/p\right).
\end{equation}
对孤立的无穷远指标 $x=nb^3$，$\phi$ 的分段定义
首先丢弃每个满足 $(nb,\mathfrak E)\ne1$ 的项。 这是
在整个完备化指标上的掩码：$(nb^3,\mathfrak E)=1$ 当且仅当 $n$
与 $b$ 都是 $\mathfrak E$-单位。 仅在这个剩余的域上，我们才使用
完全乘性，包括 $\chi_b(\lambda)^6=1$，写出
\[
 \phi(nb^3)\prod_p\chi_p(nb^3)^{j_p}
 =\chi_n(\lambda)^2\Psi(n)\Psi(b)^3
 \qquad((nb,\mathfrak E)=1).
\]
当 $n,b$ 有公共移动素数时它仍然成立，因为相应的
局部幂保持其零值，包括 $j_p=0$ 的情形。 对于
被丢弃的项，完备化系数直接定义为零；
那里不计算 $\chi_n(\lambda)$ 的任何值。 在保留域上
结合方程~\eqref{eq:reflection-infinity-coefficient} 与
\eqref{eq:reflection-cubic-sextic-conversion}，现在
把直接的 Mellin 级数与乘积 $T$ 等同起来。 对偶系数
函数 $d_H$ 保持方程~\eqref{old-eq:4.3} 中的完整源支撑；
不在它们上面放置 $S$-掩码。

固定 $h_0$，并恰在 $h_p\ne0$ 时称 $p$ 为活跃的。
方程~\eqref{eq:reflection-local-fourier} 迫使所有 $j_p\ne0$ 的
素数都是活跃的。 对每个 $h_0\bmod L$ 选取一次
代表元及既约表达式
$\lambda^2h_0/L=a_F/c_F$。 把分子与分母乘以
一个单位，当 $\lambda\mid c_F$ 时把分子正规化为 $1\pmod3$，
否则把分母正规化为 $1\pmod3$。 由于每个活跃
素数都是主的，该单位仅依赖于 $h_0$。
对 $r=\prod_{p\ {\rm active}}p$，方程~\eqref{eq:reflection-translated-theta}
随后给出
\begin{equation}\label{eq:reflection-cusp-fraction}
 z_h=\frac ac,\qquad c=c_Fr,\qquad
 a=a_Fr+\lambda^2c_F\sum_{p\mid r}(r/p)h_p.
\end{equation}
在整除 $c_F$ 的素数处，分子同余于 $a_Fr$ 且
是单位。 在 $p\mid r$ 处，它同余于
$\lambda^2c_F(r/p)h_p$，同样是单位。 因此这个分数是
既约的，即使 $h_0$ 模 $L$ 是非单位。 允许空乘积
$r=1$。 把任一提升换成另一个会使 $z_h$ 改变
$\lambda^2\mathcal O=3\mathcal O$ 中的一个元素，即 $\theta$ 的一个周期。

\emph{尖点矩阵及其自守因子的选取。}
接下来的构造对刚得到的每个既约尖点分数确定
自守因子与反射加性相位。
置 $M=\lambda^{12}L^4$。 对每个活跃的 $p$，中国剩余
定理容许把所选的非零类 $h_p\bmod p$ 提升为满足
$h_p\equiv0\pmod{M^2}$。 把 $r$ 限制到模 $M^2$ 的一个类。
方程~\eqref{eq:reflection-cusp-fraction} 表明 $a$ 此时模 $Mc_F$ 固定；
事实上其移动项被 $M^2c_F$ 整除，而 $a_Fr$ 模 $M^2$ 固定，
且由于 $c_F$ 在相差一个单位的意义下整除 $L$，$M^2$ 被 $Mc_F$ 整除。
它还表明，当 $\lambda\mid c$ 时 $a\equiv1\pmod3$，否则
$c\equiv1\pmod3$。

由以下相容的局部同余选取 $\delta'$：
\[
 \begin{array}{ll}
 \delta'\equiv a^{-1}\pmod{Mc_F}
       &\text{在整除 }c_F\text{ 的每个素数处},\\
 \delta'\equiv0\pmod M
       &\text{在整除 }M\text{ 但不整除 }c_F\text{ 的每个素数处},\\
 \delta'\equiv a^{-1}\pmod p& (p\mid r).
 \end{array}
\]
第一行指 $Mc_F$ 的所示素幂部分。
所有逆都存在，因为 $a/c$ 既约。 置
$b=(a\delta'-1)/c$ 及 $g=\bigl(\begin{smallmatrix}a&b\\c&\delta'\end{smallmatrix}\bigr)$。
则 $g\in\Gamma$，且用 $c$ 除 $a\delta'-1$ 给出
\[
 b\equiv0\pmod M\text{ 在 }c_F\text{ 的素数处},\qquad
 b\equiv-c^{-1}\pmod M\text{ 在 }M\text{ 的其他素数处}.
\]
于是 $a,b,c,\delta'$ 在整除 $L$ 的每个素数的所需幂处
有固定剩余，与类 $h_p$ 无关。
每当在剩余符号计算中出现零元素时，
可以先把 $z_h$ 平移 $\lambda^2M^2\mathcal O$，然后
给 $\delta'$ 加上组合同余模的一个倍数。
这些改变保持所有展示的同余，并避开使
某个元素为零的有限多个取值。 它们不改变平移后的
theta 函数。

有三种情形，由 $v_\lambda(c)$ 区分。
若 $3\mid c$，置 $H=I$。 则 $G=g\in\Gamma_1(3)$，而对主 $a$
与 $r$ 的主素数应用三次
互反律给出
\begin{equation}\label{eq:reflection-cusp-case-zero}
 \kappa(G)=\left(\frac{c_F}{a}\right)_3
             \left(\frac ar\right)_3.
\end{equation}
若 $v_\lambda(c)=1$，取唯一的 $u_0\in\{\lambda,-\lambda\}$
使 $u_0\equiv c\pmod3$，并置 $H=\mathsf L(u_0)$；由于
$\mathcal O/(\lambda)\simeq\mathbb F_3$，这就是两种可能。
上面的同余表明 $G=gH^{-1}\in\Gamma_1(3)$。
记 $A=a-u_0b$。 $A$ 与 $a$ 都是主的，且它们互素：
$(a,b)=1$ 且 $(a,u_0)=1$。 行列式方程给出
\[
 a(c-u_0\delta')=cA-u_0,
 \qquad bc\equiv-1\pmod a.
\]
由此，对 $a,A$ 应用三次互反律，得到
\[
 \begin{split}
 \kappa(G)
 &=\left(\frac{-u_0}{A}\right)_3
          \left(\frac aA\right)_3^{-1}\\
 &=\left(\frac{-u_0}{A}\right)_3
          \left(\frac{c/u_0}{a}\right)_3
  =\left(\frac{-u_0}{A}\right)_3
          \left(\frac{c_F/u_0}{a}\right)_3
          \left(\frac ar\right)_3.
 \end{split}
\]
对第二个等式，把 $A$ 模 $a$ 约化并使用
$(A/a)_3=(-u_0b/a)_3=(u_0/a)_3(c/a)_3^{-1}$；$-1$ 的三次符号
为一。 最后设 $(\lambda,c)=1$。 选取
$u_0=s+t\omega\equiv a\pmod3$ 且 $0\le s,t<3$，并置
$H=\bigl(\begin{smallmatrix}u_0&-1\\1&0\end{smallmatrix}\bigr)$。
这里 $c\equiv1$、$\delta'\equiv0$ 且 $b\equiv-1\pmod3$，所以
\[
 G=gH^{-1}=\begin{pmatrix}-b&a+bu_0\\-\delta'&c+\delta'u_0\end{pmatrix}
       \in\Gamma_1(3).
\]
$-b$ 与 $c$ 都是主的。 此外 $-bc=1-a\delta'\equiv1\pmod9$。
把 $\delta'=u\lambda^k\delta_0$ 分解，其中 $\delta_0$ 为主。
补充律给出 $(u\lambda^k/(1-a\delta'))_3=1$；三次
互反律给出 $(\delta_0/(1-a\delta'))_3=1$，因为互反后
分子为 $1\pmod{\delta_0}$。 于是
\[
 \left(\frac{\delta'}{-bc}\right)_3=1,
 \qquad
 \kappa(G)=\left(\frac{\delta'}{-b}\right)_3
          =\left(\frac{\delta'}c\right)_3^{-1}
          =\left(\frac ac\right)_3
          =\left(\frac a{c_F}\right)_3\left(\frac ar\right)_3.
\]
最后一个逆使用了 $a\delta'\equiv1\pmod c$。

在三种情形中我们都证明了
\begin{equation}\label{eq:reflection-multiplier}
 \kappa(G)=\kappa_F\prod_{p\mid r}\chi_p(a)^2,
 \qquad |\kappa_F|=1.
\end{equation}
当活跃的 $h_p$ 在固定扇区内变动时，因子 $\kappa_F$ 保持不变。
若 $3\mid c$，它为 $(c_F/a)_3$：单位与
$\lambda$-幂因子由 $a\bmod9$ 确定，而互反律
由 $a$ 在 $c_F$ 的素数处的剩余确定其余的固定素数因子。 在中间情形中它为
$(-u_0/A)_3((c_F/u_0)/a)_3$，由 $A=a-u_0b$ 与 $a$ 的固定
剩余以同样的方式确定。 在最后一种情形中它为 $(a/c_F)_3$，
直接由固定的分母 $c_F$ 与
$a\bmod c_F$ 确定。 因此 $a,b,c,\delta'$ 所展示的同余
仅用 $h_0$ 与 $r\bmod M^2$ 就确定了 $\kappa_F$。
它们也确定了三个函数 $\Theta_b$ 中被选中的那一个。

\emph{反射加性相位。}
记 $x=\lambda^4\mu$ 及 $D_F=\lambda^3c_F$。 在每个 $p\mid r$ 处，
方程~\eqref{eq:reflection-cusp-fraction} 给出
$a\equiv\sigma_ph_p\pmod p$。 因此
$\delta'\equiv\sigma_p^{-1}h_p^{-1}\pmod p$。 对
$D_Fr$ 的互素因子作加性中国剩余给出精确恒等式
\begin{equation}\label{eq:reflection-additive-phase}
 \begin{split}
 \check e(-\delta'\mu/c)
 &=e(-\delta'x/(D_Fr))\\
 &=\vartheta(x)\prod_{p\mid r}e(\epsilon_ph_p^{-1}x/p),
 \qquad
 \vartheta(x)=e(-\delta'_F r^{-1}x/D_F).
 \end{split}
\end{equation}
这里 $\delta'_F$ 是 $\delta'$ 模 $D_F$ 的类，而
$\vartheta$ 的公式中的 $r^{-1}$ 模 $D_F$ 取值。
例如，$p$ 处的局部分子为
$-\delta'(D_Fr/p)^{-1}=\epsilon_ph_p^{-1}$；这验证了符号
与 $\lambda$ 的幂次。 所选的同余固定了 $\delta'_F$，
因为 $D_F$ 整除相应的局部模，而扇区固定了
$r^{-1}\bmod D_F$。 因此有限多个 $D_F$ 的一个公倍数
是所有特征 $\vartheta$ 的固定模。

对式~\eqref{eq:reflection-multiplier} 取共轭贡献
$\bar\kappa_F\prod_{p\mid r}\chi_p(\sigma_p)^{-2}\chi_p(h_p)^{-2}$。
对 $j\ne0$，乘以
式~\eqref{eq:reflection-local-fourier} 使局部和变为
\[
 q_p^{-1/2}\tau_{p,j}^{-}
 \sum_{h\ne0}\chi_p(h)^{-j-2}e(\epsilon_ph^{-1}x/p).
\]
置 $y=h^{-1}$ 后，该和是次数为 $j+2$ 的 Gauss 和。
若 $j\ne4$，其值为
$q_p^{1/2}\tau_{p,j+2}^{+}\chi_p(\epsilon_px)^{-j-2}$，在
$p\mid x$ 处也如此，因为那时两边都为零。 若 $j=4$，它是
Ramanujan 和
\[
 \sum_{y\ne0}e(\epsilon_pxy/p)=-1+q_p1_{p\mid x}.
\]
对活跃的 $j=0$，同样的替换给出
\[
 -q_p^{-1}\sum_{h\ne0}\chi_p(h)^{-2}e(\epsilon_ph^{-1}x/p)
 =-q_p^{-1/2}\tau_{p,2}^{+}\chi_p(\epsilon_px)^{-2}.
\]
这些恰好是 $\omega_{p,j}B_p(x)$；非活跃的零指数
贡献 $1-q_p^{-1}$。 这证明了每个局部因子，包括
其零延拓。 $h_0$ 与活跃子集的选择至多给出
$q_L2^{\#\{p:j_p=0\}}$ 项。 拆分 $j_p=4$ 的因子
至多把它乘以 $2^{\#\{p:j_p=4\}}$，从而证明了所述的
计数。

在这些局部和之前构造的分母与固定扇区数据
仅通过 $h_0$ 与活跃根基
$r$ 依赖于局部素数集；其余活跃频率依赖恰为刚求和
的局部依赖。 我们对每个 $h_0$ 使用相同的代表元
与正规化分数，并在这些数据重现时使用相同的固定扇区选择：
尖点函数、固定加性特征与 $\kappa_F$。 一个缺席素数
与非活跃的零指数素数具有相同的活跃根基；
因此它们其余活跃的 $\sigma_p,\epsilon_p$ 也相同。
非活跃系数 $1-q_p^{-1}$ 是唯一的改变。
这证明了陈述中的同时分支相容性。
